% Abstract
\chapter*{Abstract}
\addcontentsline{toc}{chapter}{Abstract}
\begin{otherlanguage}{english}
AI coding agents such as Claude Code and Codex run shell commands on the developer's machine
themselves; a single mistake can delete files, empty a database or rewrite a remote history.
This report presents the work carried out at WePwn on \demo{}, a safety guard that intercepts
every command an agent wants to run \emph{before} it runs, predicts its effect and returns one
decision: let it through, make it reversible by first copying what it will destroy, run it
under a runtime net, or hand it to a person and refuse it when nothing can be recovered. Every
decision leaves a receipt in a hash-chained ledger.

Starting from a beta that recognised commands with regular expressions, the internship
produced an engine that parses the line and simulates its effects on a virtual view of the
file system, a rule table written as data in which every example is a test, a single policy
and a single decider, and two runtime nets: a ptrace-based system-call guard on Linux and a
WinFsp file system on Windows. The old regular-expression path no longer decides any POSIX line, and a second front end for
PowerShell shares the same engine. The
work was done one part at a time, with tests that fail on the previous code and a
922-command decision gate that forbids any decision from becoming looser without an explicit
ruling. Offensive work found and fixed vulnerabilities, among them an elevated Windows task
that ran code the agent could modify.

\medskip
\noindent\textbf{Keywords:} AI coding agents, security, shell command analysis, simulation,
recovery, ptrace, WinFsp, PowerShell, hash-chained ledger.
\end{otherlanguage}
